自主探索有机反应机理，需要分子几何、反应路径和量子化学结果共同支持同一个机理结论。基于大语言模型（LLM）的智能体可能误判三维空间关系，也可能在必要物理检验尚未满足时宣布研究完成。

本文提出 Cloudra，一个将机理推理、确定性几何操作与科学有效性检验分开的多智能体框架。按照“智能克制”的设计原则，具有拓扑感知能力的几何算子将化学意图转化为分子结构上的空间操作。以化学实体为中心的交付 DAG 将机理主张关联至归档量子化学结果，检查驻点和反应路径，并保留尚未满足的科学要求。

我们在 CloudraBench 上评估 Cloudra。该基准包含50题，覆盖基元反应、多步机理、立体化学竞争及计算方法局限。四种工作流配置共完成200次运行，Cloudra 的平均综合得分为0.780，同一模型的 Standard 基线为0.594。在 Hard 任务中，平均得分由0.689提高至0.958。40道非 Impossible 任务中，Cloudra 有24次运行获得科学自洽性满分，Standard 为7次。行为分析进一步考察显式几何测量与结构化结果关联怎样影响局部科学判断。

Cloudra 将计算执行、几何操作与物理验收检查结合起来，依据支持机理结论的计算结果及其满足的科学要求判断研究是否完成。
